\documentclass[11pt]{article}
\usepackage[a4paper,margin=25mm]{geometry}
\usepackage{amsmath,amssymb,amsthm}
\usepackage[hidelinks]{hyperref}
\newtheorem{theorem}{Theorem}
\newtheorem{lemma}{Lemma}
\title{Matrix Rings and the McCoy Property:\\A Proof of Conjecture 00000002226}
\author{GitHub: gaochengzhecpu}
\date{}
\begin{document}
\maketitle
\begin{abstract}
For a nonzero commutative ring $R$ and $n\ge1$, the matrix ring $M_n(R)$
has the McCoy property if and only if $n=1$. The case $n=1$ is McCoy's
theorem. For $n\ge2$ an explicit pair of linear polynomials shows that
$M_n(R)$ is neither right nor left McCoy. The equivalence and both
one-sided versions are verified in Lean.
\end{abstract}

\paragraph{Definitions and reading of the source.}
Rings are associative with $1$, and $A[x]$ is the polynomial ring in a
central indeterminate. Following Nielsen~\cite{nielsen}, a ring $A$ is
\emph{right McCoy} if for all nonzero $f,g\in A[x]$ with $fg=0$ there is
a nonzero $r\in A$ with $fr=0$. It is \emph{left McCoy} if under the
same hypothesis there is a nonzero $s\in A$ with $sg=0$. It is
\emph{McCoy} if it is both. This is the property in the source: the
coefficients of a zero-divisor polynomial are annihilated by one nonzero
element.

The source considers $M_n(R)$ over a commutative ring $R$; the Chinese
version states the same equivalence. We take $R$ nonzero and $n\ge1$. These are the standing conventions for a matrix
ring; without them $M_n(R)$ is the zero ring, which is McCoy vacuously
because it has no nonzero polynomials. The last section returns to this
point.

\begin{theorem}
Let $R$ be a nonzero commutative ring and $n\ge1$. Then $M_n(R)$ is
McCoy if and only if $n=1$. The same equivalence holds for the right
McCoy property alone and for the left McCoy property alone.
\end{theorem}

The theorem follows from the next two lemmas, since $M_1(R)\cong R$ and
both properties are invariant under ring isomorphism.

\begin{lemma}[McCoy \cite{mccoy}]
Every commutative ring is McCoy.
\end{lemma}
\begin{proof}
Let $f,g\in R[x]$ be nonzero with $fg=0$. Among all nonzero $h$ with
$fh=0$ choose one of least degree $m$, and write $f=\sum_i a_ix^i$ and
$h=\sum_{j\le m}b_jx^j$ with $b_m\ne0$. If $a_ih=0$ for every $i$, then
$a_ib_m=0$ for every $i$, so $b_mf=0$ with $b_m\ne0$. Otherwise let $l$
be the largest index with $a_lh\ne0$. Then
$0=fh=(a_0+\dots+a_lx^l)h$, and the coefficient of $x^{l+m}$ gives
$a_lb_m=0$. Hence $a_lh$ is nonzero of degree less than $m$, and
$f\cdot a_lh=a_l\,fh=0$. This contradicts the choice of $h$. So some
nonzero $r\in R$ satisfies $fr=rf=0$. Applying this to $gf=0$ gives a
nonzero $s$ with $sg=0$.
\end{proof}

\begin{lemma}
Let $S$ be a nonzero ring, not necessarily commutative, and $n\ge2$.
Then $M_n(S)$ is neither right McCoy nor left McCoy.
\end{lemma}
\begin{proof}
Let $E_{ab}$ be the matrix units, so $E_{ab}E_{cd}=\delta_{bc}E_{ad}$,
and let $I$ be the identity matrix. Put
\[
 f=(I-E_{22})+E_{12}\,x,\qquad g=E_{21}-E_{11}\,x .
\]
Both are nonzero because $E_{12}\ne0$ and $E_{21}\ne0$. The three
coefficients of $fg$ are
\begin{align*}
 (I-E_{22})E_{21}&=E_{21}-E_{21}=0,\\
 E_{12}E_{21}-(I-E_{22})E_{11}&=E_{11}-E_{11}+0=0,\\
 -E_{12}E_{11}&=0,
\end{align*}
so $fg=0$. Let $r\in M_n(S)$ with $fr=0$. Then $(I-E_{22})r=0$ and
$E_{12}r=0$. The first equation gives $r=E_{22}r$, and therefore
\[
 r=E_{22}r=E_{21}E_{12}r=0 .
\]
Thus no nonzero constant annihilates $f$ on the right, and $M_n(S)$ is
not right McCoy.

For the left property put
\[
 f'=E_{12}-E_{11}\,x,\qquad g'=(I-E_{22})+E_{21}\,x .
\]
Both are nonzero. The coefficients of $f'g'$ are $E_{12}(I-E_{22})=0$,
$E_{12}E_{21}-E_{11}(I-E_{22})=E_{11}-E_{11}=0$ and $-E_{11}E_{21}=0$,
so $f'g'=0$. If $sg'=0$, then $s(I-E_{22})=0$ and $sE_{21}=0$, hence
$s=sE_{22}=sE_{21}E_{12}=0$. Thus $M_n(S)$ is not left McCoy.
\end{proof}

\section*{Formal verification and scope}
The Lean definitions \texttt{RightMcCoy}, \texttt{LeftMcCoy} and
\texttt{McCoy} are stated for an arbitrary ring $A$ and Mathlib's
polynomial ring \texttt{A[X]}, whose indeterminate is central. They are
the definitions above, with \texttt{C r} the constant polynomial.

The first lemma is proved in \texttt{rightMcCoy\_of\_comm} and
\texttt{leftMcCoy\_of\_comm} from Mathlib's McCoy theorem
\begin{center}\small\texttt{Polynomial.nmem\_nonZeroDivisors\_iff}.\end{center}
The explicit ring isomorphism \texttt{oneByOne} between $M_1(R)$ and
$R$, with the two transfer lemmas named \texttt{of\_ringEquiv}, gives
\texttt{mcCoy\_one}.

\texttt{not\_rightMcCoy\_matrix} and \texttt{not\_leftMcCoy\_matrix}
prove the second lemma for matrices over any nonzero ring indexed by
any finite type with two distinct indices. They use the polynomials
displayed above, built from Mathlib's matrix units; the products, the
nonvanishing of $f$ and $g$, and the vanishing of every annihilating
constant are all proved. The main theorem is \texttt{matrix\_mcCoy\_iff}:
for a nontrivial commutative ring and $1\le n$,
\[
 \texttt{McCoy (Matrix (Fin n) (Fin n) R)}\;\leftrightarrow\; n=1 .
\]
\texttt{matrix\_rightMcCoy\_iff} and \texttt{matrix\_leftMcCoy\_iff}
are the one-sided versions. No \texttt{sorry}, \texttt{native\_decide}
or custom axiom occurs.

The hypotheses $R\ne0$ and $n\ge1$ are the only restriction.
\texttt{mcCoy\_of\_subsingleton} records why they are needed: a zero
ring is McCoy vacuously, so for $R=0$ or $n=0$ the ring $M_n(R)$ is
McCoy for a trivial reason. Under the usual convention that the
coefficient ring of a matrix ring is nonzero, the source statement is
proved in full. A reader who admits the zero ring should read the
result as: the equivalence holds exactly for nonzero $R$.

\section*{Source and reproduction}
The exact bilingual source is included byte-for-byte as
\texttt{SOURCE.md}. Its repository path is
\begin{center}\small\texttt{conjectures/00000002226.md}.\end{center}
The project pins Lean 4.19.0 and Mathlib commit
\begin{center}\small\texttt{c44e0c8ee63ca166450922a373c7409c5d26b00b}.\end{center}
From \texttt{lean/}, run \texttt{lake build}, then
\begin{center}\small\texttt{lake env lean -DwarningAsError=true Main.lean}.\end{center}
The included public-Git manifest locks all dependency revisions. The
\texttt{verification/} directory records the fresh-build logs,
theorem-axiom reports, artifact hashes, review and final PDF
inspection. No supplementary numerical computation is required.

\begin{thebibliography}{9}
\bibitem{mccoy} N. H. McCoy, Remarks on divisors of zero,
\emph{Amer. Math. Monthly} 49 (1942), 286--295.
\bibitem{nielsen} P. P. Nielsen, Semi-commutativity and the McCoy
condition, \emph{J. Algebra} 298 (2006), 134--141.
\end{thebibliography}
\end{document}
